\documentclass[tikz,border=3pt]{standalone}
\input{tikz-tcee}
% colores-solucion.tex — Colores de los gráficos propios de la solución (1D01,
% examen de España). Se cargan después de tikz-tcee.tex (paleta «dic…»).
% Los 7 primeros son los que añade el cuadernillo del examen (mismo HEX, para
% que lo mismo tenga el mismo color en el examen y en la solución; CRITERIOS §7).
\definecolor{dicbienes}{HTML}{2C6E9B}
\definecolor{dicturismo}{HTML}{8CC3E6}
\definecolor{dicrentas}{HTML}{A6A6A6}
\definecolor{dicsaldo}{HTML}{1A1A1A}
\definecolor{dicpiin}{HTML}{8C510A}
\definecolor{dicfiscal}{HTML}{4A5A6A}
\definecolor{dicvivienda}{HTML}{B07AA1}
% Nuevos en la solución:
% sector privado (gama cálida, la del gasto privado de la guía)
\definecolor{dicprivado}{HTML}{C76A1F}
% efecto bola de nieve (azul medio, distinto del pizarra de lo fiscal)
\definecolor{dicbola}{HTML}{4F86C6}
% ajuste déficit-deuda y otras variaciones (residuos): gris claro con trama
\definecolor{dicresiduo}{HTML}{B4B4B4}
% previsión del PIB (el morado del PIB, aclarado)
\colorlet{dicpibprev}{dicpib!35!white}
\pgfplotsset{
  dicsol/.style={width=13cm, height=6.4cm, scale only axis=false,
    axis lines*=left, tick align=outside, clip=false,
    ymajorgrids, grid style={tceegris!25},
    ticklabel style={font=\footnotesize}, label style={font=\footnotesize},
    xticklabel style={/pgf/number format/1000 sep=},
    yticklabel style={/pgf/number format/.cd, use comma, 1000 sep={}},
    legend style={at={(0.5,-0.15)}, anchor=north, draw=none, font=\footnotesize,
      /tikz/every even column/.append style={column sep=8pt}},
    legend cell align=left},
  leyendalinea/.style={legend image code/.code={\draw[##1] (0cm,0cm) -- (0.55cm,0cm);}},
  % leyenda de una serie de solo puntos (en un gráfico con area legend)
  leyendamarca/.style={legend image code/.code={\draw[mark repeat=2,mark phase=2,##1] plot coordinates {(0cm,0cm) (0.3cm,0cm) (0.6cm,0cm)};}},
}

\usetikzlibrary{shapes.geometric}
\begin{document}
\hyphenpenalty=10000 \exhyphenpenalty=10000
% Solución, gráfico s09. Los umbrales oficiales (infografía de la guía,
% esquema-umbrales-pdm) con la posición de España: rombo negro con el último
% dato del cuadernillo (2025, o el periodo de 3 años que acaba en 2025).
% Escala de cada fila: propia y solo indicativa (lo que cuenta es el lado del umbral).
% Capas: barras → marcas de umbral → nombres → valores del umbral (encima, sin
% fondo) → España (rombo) → valor de España (debajo, sin fondo) → leyenda.
\begin{tikzpicture}[tcee, font=\scriptsize,
  ok/.style={fill=tceeverde!70}, mal/.style={fill=tceerojo!75},
  n/.style={anchor=east, font=\scriptsize, text width=3.4cm, align=right},
  u/.style={font=\scriptsize\bfseries, anchor=south, inner sep=1pt},
  e/.style={font=\scriptsize, anchor=north, inner sep=1pt},
  esp/.style={diamond, fill=dicsaldo, draw=white, line width=0.5pt, inner sep=1.7pt}]
  % 1: barras
  \capa{1}{\fill[mal] (0,0.00) rectangle (1.2,0.30);
    \fill[ok] (1.2,0.00) rectangle (3.8,0.30);
    \fill[mal] (3.8,0.00) rectangle (5,0.30);
    \fill[mal] (0,-0.90) rectangle (2.0,-0.60);
    \fill[ok] (2.0,-0.90) rectangle (5,-0.60);
    \fill[ok] (0,-1.80) rectangle (2.5,-1.50);
    \fill[mal] (2.5,-1.80) rectangle (5,-1.50);
    \fill[ok] (0,-2.70) rectangle (2.5,-2.40);
    \fill[mal] (2.5,-2.70) rectangle (5,-2.40);
    \fill[mal] (0,-3.60) rectangle (1.6,-3.30);
    \fill[ok] (1.6,-3.60) rectangle (3.4,-3.30);
    \fill[mal] (3.4,-3.60) rectangle (5,-3.30);
    \fill[ok] (0,-4.50) rectangle (2.5,-4.20);
    \fill[mal] (2.5,-4.50) rectangle (5,-4.20);
    \fill[ok] (0,-5.40) rectangle (2.5,-5.10);
    \fill[mal] (2.5,-5.40) rectangle (5,-5.10);
    \fill[ok] (0,-6.30) rectangle (2.5,-6.00);
    \fill[mal] (2.5,-6.30) rectangle (5,-6.00);
    \fill[ok] (0,-7.20) rectangle (2.5,-6.90);
    \fill[mal] (2.5,-7.20) rectangle (5,-6.90);
    \fill[ok] (0,-8.10) rectangle (2.5,-7.80);
    \fill[mal] (2.5,-8.10) rectangle (5,-7.80);}
  % 2: marcas de los umbrales
  \capa{2}{\draw[very thick, dicinflacion] (1.2,-0.05) -- (1.2,0.35);
    \draw[very thick, dicinflacion] (3.8,-0.05) -- (3.8,0.35);
    \draw[very thick, dicinflacion] (2.0,-0.95) -- (2.0,-0.55);
    \draw[very thick, dicinflacion] (2.5,-1.85) -- (2.5,-1.45);
    \draw[very thick, dicinflacion] (2.5,-2.75) -- (2.5,-2.35);
    \draw[very thick, dicinflacion] (1.6,-3.65) -- (1.6,-3.25);
    \draw[very thick, dicinflacion] (3.4,-3.65) -- (3.4,-3.25);
    \draw[very thick, dicinflacion] (2.5,-4.55) -- (2.5,-4.15);
    \draw[very thick, dicinflacion] (2.5,-5.45) -- (2.5,-5.05);
    \draw[very thick, dicinflacion] (2.5,-6.35) -- (2.5,-5.95);
    \draw[very thick, dicinflacion] (2.5,-7.25) -- (2.5,-6.85);
    \draw[very thick, dicinflacion] (2.5,-8.15) -- (2.5,-7.75);}
  % 3: nombres de los indicadores
  \capa{3}{\node[n] at (-0.1,0.15) {Cuenta corriente (media de 3 años)};
    \node[n] at (-0.1,-0.75) {PIIN};
    \node[n] at (-0.1,-1.65) {Tasa de paro};
    \node[n] at (-0.1,-2.55) {CLU nominales (3 años, euro)};
    \node[n] at (-0.1,-3.45) {Tipo de cambio efectivo real (3 años, euro)};
    \node[n] at (-0.1,-4.35) {Precio nominal de la vivienda (anual)};
    \node[n] at (-0.1,-5.25) {Deuda de los hogares};
    \node[n] at (-0.1,-6.15) {Deuda de las empresas};
    \node[n] at (-0.1,-7.05) {Flujo de crédito (hogares / empresas)};
    \node[n] at (-0.1,-7.95) {Déficit / deuda públicos (Tratado)};}
  % 4: valores del umbral (encima de la marca)
  \capa{4}{\node[u] at (1.2,0.36) {$-4$};
    \node[u] at (3.8,0.36) {$+6$};
    \node[u] at (2.0,-0.54) {$-35$};
    \node[u] at (2.5,-1.44) {10};
    \node[u] at (2.5,-2.34) {$+9$};
    \node[u] at (1.6,-3.24) {$-3$};
    \node[u] at (3.4,-3.24) {$+3$};
    \node[u] at (2.5,-4.14) {$+9$};
    \node[u] at (2.5,-5.04) {55};
    \node[u] at (2.5,-5.94) {85};
    \node[u] at (2.5,-6.84) {14 / 13};
    \node[u] at (2.5,-7.74) {3 / 60};}
  % 5: España (rombo)
  \capa{5}{\node[esp] at (3.05,0.15) {};
    \node[esp] at (1.69,-0.75) {};
    \node[esp] at (2.65,-1.65) {};
    \node[esp] at (3.24,-2.55) {};
    \node[esp] at (2.23,-3.45) {};
    \node[esp] at (3.06,-4.35) {};
    \node[esp] at (2.02,-5.25) {};
    \node[esp] at (1.58,-6.15) {};
    \node[esp] at (1.1,-7.05) {};
    \node[esp] at (2.32,-7.95) {};
    \node[esp] at (4.12,-7.95) {};}
  % 6: valor de España (debajo de la barra)
  \capa{6}{\node[e] at (3.05,-0.05) {3,1};
    \node[e] at (1.69,-0.95) {$-42{,}7$};
    \node[e] at (2.65,-1.85) {10,5};
    \node[e] at (3.24,-2.75) {$\approx$13,9*};
    \node[e] at (2.23,-3.65) {$-0{,}9$};
    \node[e] at (3.06,-4.55) {12,7};
    \node[e] at (2.02,-5.45) {42,9};
    \node[e] at (1.58,-6.35) {62,0};
    \node[e] at (1.1,-7.25) {4,5 / 4,2*};
    \node[e] at (2.32,-8.15) {2,4};
    \node[e] at (4.12,-8.15) {100,5};}
  % 7: leyenda y nota
  \capa{7}{\node[esp] at (-3.4,-8.85) {};
    \node[anchor=west, font=\scriptsize] at (-3.25,-8.85) {España (último dato del cuadernillo)};
    \draw[very thick, dicinflacion] (1.3,-9.00) -- (1.3,-8.70);
    \node[anchor=west, font=\scriptsize] at (1.4,-8.85) {umbral};
    \fill[ok] (2.45,-8.95) rectangle (2.75,-8.75); \node[anchor=west, font=\scriptsize] at (2.8,-8.85) {sin alerta};
    \fill[mal] (4.15,-8.95) rectangle (4.45,-8.75); \node[anchor=west, font=\scriptsize] at (4.5,-8.85) {alerta};
    \node[anchor=north west, text width=8.9cm, font=\scriptsize\itshape, text=tceegris] at (-3.5,-9.10)
    {En \% (del PIB, salvo el paro, las variaciones y el flujo de crédito, sobre la deuda del año anterior). España: 2025 o 2023-2025 (cuadros 3, 4, 6, 7 y 9). (*) Aproximaciones: CLU por persona (el umbral es por hora); flujo de crédito aproximado por la variación de los préstamos a hogares y empresas. La posición dentro de cada barra es indicativa.};}
\end{tikzpicture}
\end{document}
